\section*{Введение}
\addcontentsline{toc}{section}{Введение}

\textbf{Актуальность темы.} В условиях стремительного роста популярности аниме-культуры во всём мире и увеличения объёма доступного контента пользователи сталкиваются с проблемой информационной перегрузки. Ежегодно выпускаются сотни новых тайтлов, что затрудняет поиск релевантных произведений и формирование персонализированной зрительской траектории. Традиционные методы навигации по каталогам (по жанрам, году, рейтингу) не учитывают индивидуальные предпочтения пользователя и его историю взаимодействий. Разработка специализированных веб-приложений, интегрирующих современные алгоритмы рекомендательных систем и средства социального взаимодействия (комментарии, избранное), является необходимым условием для повышения качества пользовательского опыта и удержания аудитории аниме-сообществ.

\textbf{Степень разработанности проблемы.} Задача построения персонализированных рекомендательных систем классифицируется в области машинного обучения как задача фильтрации контента. Основы её решения заложены в работах Решника (коллаборативная фильтрация) и Паццани (контентно-ориентированные подходы). В области веб-ориентированных аниме-платформ существенный вклад внесли проекты Shikimori, MyAnimeList и AniDB. Несмотря на наличие коммерческих и открытых решений, создание гибких микросервисных архитектур с гибридными рекомендациями (content-based + collaborative), адаптированных под специфику аниме-контента (жанры, студии, режиссёры, оценки), остаётся актуальной и не до конца решённой практической задачей.

\textbf{Объект исследования} — процессы организации и управления пользовательским контентом в аниме-энциклопедиях, включая каталогизацию произведений, сбор пользовательских оценок и формирование рекомендаций.

\textbf{Предмет исследования} — программные средства и алгоритмические решения для автоматизации построения персонализированных рекомендаций, управления пользовательскими сессиями (аутентификация, избранное, комментарии) и контент-наполнением в рамках веб-ориентированной платформы на основе микросервисной архитектуры.

\textbf{Целью данной работы} является проектирование и реализация клиент-серверного веб-приложения (аниме-энциклопедии), предназначенного для предоставления информации о произведениях, ведения пользовательских активностей (избранное, комментарии) и формирования персонализированных рекомендаций за счёт гибридного подхода, сочетающего коллаборативную и контентную фильтрацию.